\section{Related Work}
\label{sec:related-work}

Recent surveys organize MoE inference optimization across model-, system-, and hardware-level techniques~\cite{liu2026survey}. Argine is closest to work on memory-constrained inference, expert offloading, caching, and prediction, but differs in experimental emphasis: rather than proposing a single offloading policy, it characterizes the allocation, retention, metric, and pressure-regime assumptions on which such policies depend.

\subsection{Sparse MoE architectures and distributed inference}

Modern sparse MoE systems build on sparsely gated conditional computation, in which a router activates only a small subset of a much larger expert pool~\cite{shazeer2017sparsely}. GShard demonstrated large-scale conditional computation with automatic sharding~\cite{lepikhin2020gshard}, while Switch Transformers simplified sparse routing and scaled MoE models to trillion-parameter regimes~\cite{fedus2022switch}. These works establish the architectural motivation for large expert pools whose total parameter footprint can greatly exceed the parameters active for one token.

Systems such as DeepSpeed-MoE and Tutel primarily address efficient MoE execution at data-center scale through compression, parallelism, communication optimization, and adaptive distributed execution~\cite{rajbhandari2022deepspeedmoe,hwang2023tutel}. Argine studies a complementary regime: a single consumer CPU with a hard DRAM budget, where expert weights are streamed from NVMe rather than distributed across an accelerator cluster. The resulting bottleneck is not all-to-all communication between expert-parallel workers, but which expert weights occupy a small resident pool and how non-resident demand becomes exposed storage time.

\subsection{Memory-constrained LLM inference}

Memory pressure has also motivated dense-model inference systems that span GPU memory, host memory, and storage. FlexGen searches for tensor placement and access schedules across GPU, CPU, and disk to maximize throughput under limited accelerator memory~\cite{sheng2023flexgen}. PagedAttention and vLLM instead address dynamic KV-cache allocation using paging-inspired memory management~\cite{kwon2023pagedattention}. These systems show that memory placement and movement policy can be first-order inference concerns even outside MoE.

Argine differs in both granularity and objective. Its managed objects are sparsely requested expert-weight slices, while non-expert weights remain outside the expert-memory budget. The central question is therefore not a static placement of full layers or KV blocks, but the interaction between router-driven demand, layer allocation, expert retention, and asynchronous storage reads at batch-one generation.

\subsection{Expert offloading and heterogeneous execution}

A growing line of work exploits MoE sparsity directly. Eliseev and Mazur showed that expert-aware offloading can make models such as Mixtral practical on limited consumer hardware~\cite{eliseev2023offloading}. EdgeMoE partitions MoE models across an edge-device storage hierarchy and combines expert prediction with bit-width adaptation~\cite{yi2023edgemoe}. MoE-Infinity uses request-level activation traces to drive a sparsity-aware expert cache, prefetching, and replacement on personal machines~\cite{xue2025moeinfinity}. These systems are particularly close to Argine in recognizing that only a subset of experts should occupy fast memory at a time.

Other systems reduce transfer cost by changing where or how missed experts execute. Fiddler orchestrates CPU and GPU execution to reduce CPU--GPU movement when GPU memory is insufficient~\cite{kamahori2025fiddler}. AdapMoE combines adaptive expert gating and expert management for efficient inference~\cite{zhong2024adapmoe}. Klotski uses expert-aware multi-batch pipelining and correlation-aware prefetching to overlap expert I/O with computation~\cite{fang2025klotski}. MoE-APEX further uses adaptive precision for cache-miss experts together with prefetching and cost-aware caching~\cite{tang2026moeapex}.

Argine does not compete with these systems as an end-to-end GPU offloading engine. Its contribution is orthogonal: it measures how much opportunity is attributable to layer allocation versus within-layer selection, validates replay semantics against the runtime, separates reuse metrics from elapsed-time prediction, and tests whether apparently natural temporal and normalized-pressure explanations survive direct experiments. These questions matter to offloading systems because a high hit rate, an adaptive partition, or a normalized memory ratio is useful only insofar as it predicts or reduces exposed latency.

\subsection{Expert prediction, prefetching, and cache-aware routing}

Expert prediction attempts to hide transfer latency by learning or exploiting structure in router decisions. Cache-Conditional Experts modifies routing itself to favor cached experts on memory-constrained mobile devices~\cite{skliar2024cacheconditional}. Local-routing-consistency analysis shows that the amount of exploitable temporal locality varies substantially across MoE architectures, cautioning against assuming that all models are equally cacheable~\cite{liang2025routingconsistency}. Pre-attention expert prediction uses same-layer pre-attention activations to forecast routing early enough for prefetching~\cite{zhu2025preattention}. Fate uses cross-layer gate information and a shallow-favoring cache policy for edge inference~\cite{fang2026fate}, while DoMoE exploits domain locality to narrow semantic expert prediction~\cite{mu2026domoe}.

These approaches motivate the large future-informed selection gap observed in Section~\ref{sec:allocation-selection}: there is substantial value in predicting \emph{which} experts will be needed, not only in deciding how much memory each layer receives. Argine's Belady-style reference is intentionally not a prediction algorithm. Following the classical role of optimal offline replacement~\cite{belady1966replacement}, it serves as a future-informed reference against which the residual retention opportunity can be measured. The paper therefore does not claim that the measured headroom is online-achievable.

\subsection{Dynamic allocation and online serving}

Dynamic resource allocation appears at multiple levels of MoE systems. Distributed serving can reassign accelerator resources as expert popularity changes; recent theoretical work formalizes this as an online MoE-serving problem with reconfiguration cost~\cite{huang2026moeserving}. Argine considers a different resource: bytes of a fixed local expert-memory pool. Nevertheless, both settings expose the same broad tension between reacting to a changing demand distribution and paying the cost of changing the allocation.

The transition results in Section~\ref{sec:workload-drift} add an empirical qualification to this literature. Strong routing drift is observable, but it does not follow that a reactive controller is useful: the tested \texttt{adaptive-e64} temporal mechanism is outperformed by a timing-free mean allocation in the evaluated transitions. This distinguishes the existence of non-stationarity from the value of a particular adaptation rule and motivates future controllers that evaluate movement in terms of predicted time savings rather than change detection alone.

\subsection{Position of this work}

The closest prior systems generally optimize one or more of expert prediction, caching, prefetching, quantization, heterogeneous execution, or I/O overlap. Argine instead treats the expert store as an experimental object and asks which abstractions remain valid under a hard memory ceiling. Its main differences are: (i) CPU-only NVMe-to-DRAM expert streaming rather than host-to-GPU expert offload; (ii) explicit separation of layer allocation from within-layer retention; (iii) event-atomic replay validated against runtime behavior; (iv) empirical separation of runtime hit semantics, expert-level reuse, and time cost; (v) preregistered negative tests of temporal adaptation and cross-model pressure explanations; and (vi) extreme-pressure measurements below the one-token expert footprint.

Accordingly, the paper should be read as a characterization and systems-analysis complement to expert-offloading engines, not as a replacement for them. Its results identify assumptions that future caching and prefetching mechanisms can use as design constraints and, equally importantly, assumptions they should not take for granted.
